\section{Classical vertices after stationary elimination}
\label{sec:ex-classical-stationary-vertices}

We now construct the cubic and quartic coefficient pictures of the
stationary action. Two estimates must be kept distinct. The fine action
is controlled by its vanishing order in momentum; the internal contraction
in the quartic term is controlled in the weighted fluctuation norm. A
coordinatewise estimate for that contraction would introduce an unwanted
factor equal to the number of fine variables in a block.

\subsection{The weighted vertex hypotheses}
Write $\mu=L^{-1}$, where $L$ is the block side. All momenta in this
section are coarse momenta. Let $\Phi(p)$ be the linear fine-field response
to a coarse field. In block coordinates the full tangent vector is
\[
        \mathcal L_p h=(h,\Phi(p)h,0).
\]
The last component is the gauge coordinate. An internal vector is written
$\iota v=(0,v)$, with the internal space including the fine field and the
auxiliary scalar variables. Let $\|\cdot\|_p$ be its weighted fluctuation
norm. The corresponding bilinear dual norm is
\[
 \|s\|_{p,*}=\sup_{\|v\|_{-p}\leq1}|v^Ts|.
\]
The sign in this definition is dictated by momentum conservation.

Fix $0<\eta\leq1$ and $D\geq0$. The quantitative bridge hypothesis is
holomorphy and periodicity together with
\begin{equation}\label{eq:ex-vertex-cruise}
 \sup_z|\nabla^{j}\Phi(p)e_a(z)|\leq D\mu^{1+j},
 \qquad 0\leq j\leq2,\quad p\in\operatorname{strip}(\eta).
\end{equation}
Here $e_a$ is a unit color-direction coordinate and $\nabla^j$ is any
ordered string of $j$ fine differences, with the coarse Fourier phase
included at a block crossing. Choose $D_*=1+D+5e^4$. The same bound with $D$ replaced by
$D_*$ applies to the fine component of $\mathcal L_pe_a$.
For one internal vector of norm at most one, the weighted discrete
Fourier estimate supplies the analogous bounds with one unit of
integrability charge, through difference order one. More explicitly,
if $H_b$ denotes the normalized block norm $L^{2/b}$ for $b>0$ and
$L^\infty$ for $b=0$, the required estimates are
\[
 H_0(\nabla^j\mathcal L_pe_a)\leq D_*\mu^{1+j},\qquad
 H_1(\nabla^j\iota v)\leq D_*\mu^{1+j}\quad(j\leq1).
\]
Only the fine-field component occurs in these two inequalities.
The weighted Fourier estimate and the block penalty estimate below are
explicit inputs here; an ordinary bound for chart derivatives at fixed
$L$ would not imply either uniform estimate.

An $n$-linear vertex has price $C_n$ if its absolute value is at most
$C_n$ on smooth coordinate legs, and also with the first leg replaced by
$\iota v$, $\|v\|_p\leq1$, and all other legs smooth. These are
\emph{unnormalized} derivatives. Their coefficient pictures therefore
carry a factor $1/n!$.

For the block penalty we use the following precise receipt. For $n=3,4$,
its real $n$th derivative in each of the four coarse directions is bounded
by $B_n$, on the stencil fields induced by these legs, either all smooth
or with one internal leg. The constants $B_n$ depend on the fixed group,
patch enlargement, penalty coefficient and $D$, but not on $L$ or the
incoming interaction budget. The stencil field bounds are to hold for
both real and imaginary parts. Complex multilinearity then gives price
\begin{equation}\label{eq:ex-vertex-penalty-price}
                    N_n=2^n\,4B_n.
\end{equation}
Indeed, expand each complex leg into its real and imaginary parts.
There are $2^n$ terms, each with a coefficient of modulus one; sum the
four directional penalties. This deduction does not replace the required
real weighted stencil estimate.

\subsection{Why the residual vertices gain a scale factor}
The coefficient pictures proved in
Section~\ref{sec:ex-soft-curvature-receipts} split the fine classical
action into its Wilson part and its residual part. At degree three the
Wilson picture vanishes to order one and the residual picture to order
two; at degree four their respective orders are zero and one.
Lemma~\ref{lem:ex-soft-factorization} writes each as a finite sum of
products of one-link difference factors and analytic coefficient arrays.
The factors may be assigned to legs other than the first. Thus the leg
that will later be tested in the internal norm carries no differences.
For the four cases, the number of legs $n$ and the number of differences
$q$ satisfy
\[
\begin{array}{c|cc|c}
 &n&q&n+q-4\\ \hline
\text{Wilson cubic}&3&1&0\\
\text{Wilson quartic}&4&0&0\\
\text{residual cubic}&3&2&1\\
\text{residual quartic}&4&1&1.
\end{array}
\]

To see the significance of the last column, consider one term with
$j_i$ differences on leg $i$. Equation~\eqref{eq:ex-vertex-cruise}
contributes $\prod_i\mu^{1+j_i}=\mu^{n+q}$. Summation over one fine
block contributes $L^4=\mu^{-4}$. The result is
$\mu^{n+q-4}$. This count remains valid with one internal leg: use its
normalized $L^2$ bound and the supremum bounds on the other legs, then
H\"older on the normalized block measure. The displaced arguments and
analytic coefficient sums are bounded by the absolute exponentially
weighted coefficient sum furnished by the difference factorization.
Fourier phases are bounded on the chosen strip; these factors and the
finite color and difference labels are included in the constant. The
first leg has no differences, so its weaker derivative allowance is
sufficient. This proves both parts of the vertex price.

Choose the coarse width $\eta$ before fixing the incoming path width
$\tau$ and budget $A$. The Wilson arrays and the penalty prices then give
constants $a,b\geq0$ independent of $\tau,A,L$. For each $\tau,A$ the
residual arrays give a fine width $\sigma'>0$ and constants $\ell,r\geq0$.
Provided $\mu\eta\leq\sigma'$, addition of the three contributions gives
holomorphic periodic vertices $T_3,T_4$ with prices
\begin{equation}\label{eq:ex-classical-vertex-prices}
                    C_3=a+\mu\ell,\qquad C_4=b+\mu r.
\end{equation}
These vertices represent the actual block action on every torus in the
chosen family. To verify this, insert Fourier waves into each local
translation sum. Summing the coarse anchor gives zero unless the total
momentum is zero, and otherwise gives the coarse volume. The remaining
fine-block sum is exactly the array just estimated. The germ identity
between the block action and the sum of the fine action and block
penalty identifies these derivatives for the actual tags. No change of
stationary branch has been made.

\subsection{The stationary contraction formula}
Let $\mathcal A(x,y)$ be this real analytic block action, with coarse
variable $x$ and internal variable $y$. Suppose
$D_y\mathcal A(x,y(x))=0$, $y(0)=0$, and
$M=D_y^2\mathcal A(0,0)$ is positive definite. Put
$E(x)=\mathcal A(x,y(x))$ and $\mathcal Lh=(h,Dy(0)h)$.
Differentiating stationarity once gives
\begin{equation}\label{eq:ex-stationary-orthogonality}
             D^2\mathcal A(0)[\mathcal Lh,\iota v]=0.
\end{equation}
Define the internal covector
\[
 s(h,k)[v]=D^3\mathcal A(0)[\iota v,\mathcal Lh,\mathcal Lk].
\]
A second differentiation of stationarity gives
\[
                 D^2y(0)[h,k]=-M^{-1}s(h,k).
\]
The resulting third and fourth derivatives are
\begin{align}
 D^3E(0)[h_1,h_2,h_3]
   &=D^3\mathcal A(0)[\mathcal Lh_1,\mathcal Lh_2,\mathcal Lh_3],
       \label{eq:ex-stationary-third}\\
 D^4E(0)[h_1,h_2,h_3,h_4]
   &=D^4\mathcal A(0)[\mathcal Lh_1,\mathcal Lh_2,
                       \mathcal Lh_3,\mathcal Lh_4]\notag\\
   &\quad-\sum_{\{\{i,j\},\{k,l\}\}}
                      s(h_i,h_j)^TM^{-1}s(h_k,h_l),
       \label{eq:ex-stationary-fourth}
\end{align}
where the sum runs over the three unordered pairings of four indices.

For completeness, these identities can be checked without suppressing
any chain-rule terms. Along a ray, write
$y(th)=tDy(0)h+t^2u/2+t^3v/6+O(t^4)$.
The terms involving $D_y\mathcal A(0)$ vanish by stationarity. The
mixed Hessian terms involving $\mathcal Lh$ vanish by
\eqref{eq:ex-stationary-orthogonality}. The cubic coefficient is therefore
$D^3\mathcal A(0)[\mathcal Lh]^3/6$. The fourth-order coefficient is
\[
 \frac{D^4\mathcal A(0)[\mathcal Lh]^4}{24}
       +\frac{s(h,h)[u]}4+\frac{u^TMu}{8}
 =\frac{D^4\mathcal A(0)[\mathcal Lh]^4}{24}
       -\frac{s(h,h)^TM^{-1}s(h,h)}8.
\]
Here $u=-M^{-1}s(h,h)$. Polarization of these symmetric multilinear
forms proves \eqref{eq:ex-stationary-third}--\eqref{eq:ex-stationary-fourth},
including the coefficient and sign of all three pairings.

\subsection{Common analytic pictures and their prices}
Assume that the inverse of the same internal Hessian has a periodic
holomorphic Fourier symbol $V(p)$, with the weighted operator estimate
\begin{equation}\label{eq:ex-exchange-inverse}
                 \|V(p)s\|_p\leq F\|s\|_{p,*}
\end{equation}
on the strip of width $\eta$, where $F\geq0$ is fixed before the
incoming width and budget. This is an input about the actual inverse,
not an arbitrary comparison matrix.

With smooth legs at momenta $p,q$, let
\[
 s_{ab}(p,q)[v]
   =T_3(-(p+q),p,q)[\iota v,\mathcal L_pe_a,\mathcal L_qe_b].
\]
By the one-internal-leg estimate,
$\|s_{ab}(p,q)\|_{p+q,*}\leq C_3$. Consequently every contraction
in \eqref{eq:ex-stationary-fourth} has modulus at most $C_3FC_3$.
This is duality in the weighted norm; no sum of coordinatewise bounds
and no dimension factor is required. The pairing uses opposite momenta
on its two internal legs.

For $n$ external legs, take the first momentum to be minus the sum of
the other $n-1$. On the independent strip of width $\eta/(2n)$ all
external momenta and all paired internal momenta lie in the original
strip. Products, contractions and sums of the holomorphic arrays are
therefore holomorphic there. Periodicity follows from that of the
vertices, bridge and inverse. Equations
\eqref{eq:ex-stationary-third}--\eqref{eq:ex-stationary-fourth}, divided
by the factorials, define common coefficient pictures $P_3,P_4$ with
\begin{equation}\label{eq:ex-classical-picture-prices}
       \|P_3\|_{\eta/6}\leq\frac{C_3}{6},\qquad
       \|P_4\|_{\eta/8}\leq\frac{C_4+3C_3FC_3}{24}.
\end{equation}
They record the stationary derivatives exactly on every torus: the
vertex identities already have the coarse-volume normalization, and
Fourier inversion of the internal Hessian makes its contraction precisely
$V$ at the paired momentum. Thus the contraction formula preserves the
same momentum-conservation rule and the same volume factor.

Finally choose $L$ sufficiently large that $\mu\ell,\mu r\leq1$.
Both pictures are then bounded on their common strip by
\[
 M_*=1+\frac{a+1}{6}
          +\frac{(b+1)+3(a+1)F(a+1)}{24}.
\]
The width and $M_*$ are chosen before $\tau,A$; only the lower threshold
for $L$ depends on the incoming path. This order of choices is what
allows the stationary output construction to be iterated.
The weighted bridge, weighted penalty and common inverse estimates
remain the precise supplier interfaces of this section.
